% ==============================================================================
% T0 Theory / FFGFT: Document 378 (Chapter Content, EN)
% Title: Peer Review under Present Conditions
% Author: Johann Pascher
% Date: 28 September 2026
% References: Doc. 190, 272, 274, 342, 351, 361, 365, 372, 377;
%   F O R M / CVP v1.0 (Vossen ed.), CIL (Vossen/Pascher), CCR (Vossen),
%   Universal Constraint / FER (Haskins), SIM2 (Krueger).
% ==============================================================================

\providecommand{\BM}{\textbf{[B]}}
\providecommand{\KM}{\textbf{[K]}}
\providecommand{\SM}{\textbf{[S]}}
\providecommand{\XM}{\textbf{[X]}}
\providecommand{\QM}{\textbf{[Q]}}

\chapter*{Peer Review under Present Conditions}
\addcontentsline{toc}{chapter}{Peer Review under Present Conditions}

{\small\noindent\textbf{Abstract.}\quad\ignorespaces
Classical peer review assumes that a referee can survey a paper in a few hours, that
author and referee belong to an institution, and that texts are written by people.
None of these assumptions holds in general any more: corpora run to hundreds of
documents, independent researchers have no access to the usual submission routes, and
language models produce plausible texts, scripts and apparent review reports in any
quantity. This document collects the review tools that have already arisen in the
exchange of the \emph{Information Physics Institute} (IPI), in Stefaan Vossen's Dot
Theory governance and in the FFGFT corpus's own record-keeping, and asks which of them
can be used for a review framework suited to present conditions. The result is a
proposal of ten building blocks: a declared constitution before review, an entry path
provided by the author, testable objects instead of texts, predictions fixed in
advance, separation of roles, a continuing register, rules for the use of language
models and for simplifications, localised instead of global verdicts, and the explicit
acknowledgement that there is no universal checker. The document contains no physical
claims; it is a methodological inventory with the character of a proposal.
}\medskip

\section*{Legend}
\addcontentsline{toc}{section}{Legend}

\medskip
\begin{tabular}{@{}lp{11cm}@{}}
\toprule
\textbf{Label} & \textbf{Meaning} \\
\midrule
{[K]}, {[B]}, {[S]}, {[Q]}, {[X]} & FFGFT status markers: numerically confirmed, algebraically proved,
  structurally plausible with open derivation, taken from an external source, recognised as false \\
{[SETZUNG]} & explicitly posited foundational assumption, not derived \\
IPI & Information Physics Institute, mailing list and working groups \\
F O R M & \emph{Framework -- Operator -- Residual -- Model}, the grammar of the
  \emph{Constitutional Verification Protocol} (CVP) \\
CIL & \emph{Constitutional Interpretation Layer} (Vossen/Pascher) \\
CCR & \emph{Canonical Constitutional Record} of \emph{Constitutional Physics} (Vossen) \\
FER & analytical method of Diana Haskins (Universal Constraint) \\
OAP, FAH & admissibility record and \emph{Framework Admissibility History} of Dot Theory \\
O0, O2 & originator declaration and secondary rendering (Doc.~272) \\
SHA-256 & cryptographic hash establishing the identity of a file object \\
\bottomrule
\end{tabular}

% ============================================================================
\section{Starting point}
% ============================================================================

\subsection{What the classical procedure assumes}

Journal peer review is built for the article: a finished paper of ten to thirty pages,
one referee from the same field, a few hours of reading, a verdict in three grades
(accept, revise, reject). It rests on four tacit assumptions:
\begin{enumerate}
\item The work can be surveyed in one sitting.
\item The referee shares with the author a common stock of methods against which
  deviations stand out.
\item Institutional affiliation acts as a prior filter.
\item A text is the product of human work, so its length is a rough measure of the effort
  behind it.
\end{enumerate}

\subsection{What has changed}

\paragraph{Scale.} A framework such as FFGFT now comprises more than three hundred and
seventy documents with their own status markers, a correction register (Doc.~190) and
cross-references between geometry, Galois theory and particle physics. Nobody takes that
in within a few hours. This is not a special case: every serious theory reaches this
point, and established programmes do not open up in an afternoon either.

\paragraph{Access.} Independent researchers have no institutional address; indexing
services often list only what is already listed through publisher platforms.
Repositories (GitHub) and archives with DOIs (Zenodo) are public and citable, but the
review system hardly registers them.

\paragraph{Language models.} Text, code and apparent review reports can now be produced
at negligible cost. The length of a text no longer measures the work behind it. In the
IPI exchange there have been cases in which a language-model prompt was submitted as
``source code'', in which a check block printed its own constants and then ``PASS''
beneath them, and in which a second language model with shared context reproduced an
earlier analysis verbatim instead of checking it independently.

\paragraph{Pattern search.} Any pattern can be read out of enough numbers. The image of
the magic lamp that one rubs until something appears describes a real misuse of computing
power. It does not describe every piece of work that contains numbers. The difference lies
in the order: in pattern search the measured values stand at the beginning, in a
derivation they stand at the end as the test. In FFGFT the geometry $T^4/\mathbb{Z}_3$ is
fixed as a {[SETZUNG]}; from it follow algebraically $\mathbb{Z}_3$, $\mathrm{GF}(27)$,
$D_4$ and triality; only from these come numbers, and only afterwards are they compared
with measurements (Doc.~365). An outsider sees at first glance only the number, not the
order, and suspects the lamp.

\paragraph{Simplification and prejudice.} A short entry path is possible and exists in
several forms. It is, however, often read as faulty in itself, because the reader takes
the simplification for the claim. Whoever comes expecting that work outside the
institutions cannot be right finds the confirmation in every simplification. The burden
therefore lies on both sides: on the author, who must build the entry path so that the
order is visible, and on the reviewer, who must consult the full version before declaring
an error.

\subsection{The question}

The question is not whether review is necessary but what it can look like under these
conditions: who reviews what, on which object, in which order, with what binding force,
and who retains interpretive authority over what.

% ============================================================================
\section{Inventory: what already exists}
% ============================================================================

The following tools were not designed for a review procedure. They arose in the work
itself, often from a concrete dispute. That is precisely why they are tested.

\subsection{The record-keeping of the FFGFT corpus}

\begin{itemize}
\item \textbf{Status markers.} Every statement carries a status: {[K]}, {[B]}, {[S]}, {[Q]},
  {[X]}, {[SETZUNG]}. The reader sees at once what is proved, what is computed, what is open
  and what is taken over. The markers existed before the CCR correspondence; later
  terminology describes what they do (Doc.~351, methodological note).
\item \textbf{Correction register.} Doc.~190 records corrections to earlier documents and the
  closing of points held open, continuously; earlier versions are frozen as an archive.
  Corrections are not worked in silently; they are recorded.
\item \textbf{Check scripts.} Every recent document comes with a script containing
  assertions that prints \texttt{PASS} or \texttt{FAIL} for each statement. Many of these
  scripts are null tests: if chance yields the same result, it is not counted.
\item \textbf{Negative results as sections.} Doc.~342 explicitly grants Galois products with
  $k\ge5$ no status from product search; Doc.~361 has a section of its own, ``What does not
  match''; Doc.~276 records a weak result as an honest negative.
\item \textbf{Guard-rail against exponent numerology.} Every quantity can be written as
  $\xi^N$; a measured scale inside an exponent is not a prediction (Doc.~190, P35).
\item \textbf{Comparator is not input.} Measured values (PDG) serve for comparison, not for
  derivation. This separation has been taken into F O R M as Illustration~001.
\end{itemize}

\subsection{The governance of Dot Theory}

Dot Theory (Vossen) is by its own account not a competing physical model but a
representational meta-programme. Its governance layer is set out in Doc.~272:
\emph{Lexicon} (terms), \emph{OAP} (admissibility), \emph{Matrix} (relations between
frameworks), \emph{FAH} (how admissibility changes). Two distinctions carry it:
\begin{itemize}
\item \emph{Preservation of meaning-making} versus \emph{allocation of meaning}: recording
  how a result was understood is admissible; deciding who is right is not.
\item Originator declaration (O0) versus secondary rendering (O2): a third party's reading
  of a framework is marked as such; the right of correction stays with the originator.
  Records are permanent and contestable; later entries may refine, dispute or supersede,
  but the earlier rendering remains part of the history.
\end{itemize}

\subsection{The Constitutional Interpretation Layer}

The CIL is co-authored by Stefaan Vossen and the author; the co-authorship is limited to
the CIL specification and does not extend to Dot Theory or the CCR. Its subject is not
the correctness of scientific statements but the entitlement to draw, inherit and pass on
conclusions. Explicitly excluded are scientific correctness, empirical validity,
mathematical truth and ontological commitment. From the dialogue between the two authors
came the principles
\begin{itemize}
\item 3 -- \emph{Asymptotic Humility},
\item 6 -- \emph{Boundary Legitimacy / Evidential Admissibility},
\item 7 -- \emph{Recoverability / Reproducibility},
\item 8 -- \emph{Asserted Constitutional Properties}: an asserted property (determinism,
  self-management, derivational status) remains an authorial assertion until it is
  independently recoverable.
\end{itemize}
Further principles of version v1.1 name the vocabulary trap (adopting the vocabulary is no
evidence of adopting the method; similarity of result or terminology is not a bridge) and
\emph{constitutional drift} (undeclared extension of scope). The author added the condition
that scope must be declared at the point of derivation and cannot be asserted
retrospectively in a dispute; this turns the principle into a checkable test.

\subsection{F O R M and the Constitutional Verification Protocol}

F O R M (ed.\ Vossen; authors Haskins, Leavitt, Pascher, Quílez Zamora, Vossen; v1.0, fixed
by SHA-256 \texttt{bd2a1f2d\ldots e8e90c0e}) divides every review into \emph{Framework}
(the declared distinction space), \emph{Operator} (the operation carried out),
\emph{Residual} (what survives the operation) and \emph{Model} (the recoverable
representation). Filling the R position with \emph{Residual} rather than \emph{Reality}
goes back to a proposal of the author: \emph{Reality} would have introduced an ontological
commitment that the protocol is meant to avoid.

The argument against a universal checker comes from the same correspondence. A proof
checker such as Lean works because axioms, type system and admissible operations were
agreed beforehand; the checker comes second. The law works the same way: an agreed,
contestable procedure precedes the ruling. Programming languages lie in between and make
the point most sharply: they are strongly normalised, and yet different libraries return
different results for the same integral, and $0.1+0.2=0.3$ is false in IEEE~754
arithmetic and true in exact rational arithmetic. If the result depends on the frame even
there, then for heterogeneous scientific frameworks a single checker cannot be the goal,
only a procedure that discloses the normalisation in each case and makes it recoverable.

\subsection{FER}

Diana Haskins' analytical method FER asks of every term what earned it its status, and
does not let terms become stronger because they fit together nicely. An application of
FER in the IPI exchange of 24--25~September 2026 produced a recursive review sequence:
declared frame $\to$ operation $\to$ distinctions obtained $\to$ representation $\to$ next
operation $\to$ recovery test $\to$ retained, lost or newly exposed distinction $\to$
revised representation. Three cases are kept apart: a distinction was available and is no
longer recoverable; it was never represented; it was not recognised, but the information
was present and a later operation exposes it. The representation is thus not the end of a
review but the next object of review.

\subsection{Practice in audits}

Several review procedures of recent months have produced rules that hold beyond their
occasion:
\begin{itemize}
\item \textbf{Freezing and hashing.} The author freezes the object and states its SHA-256
  hash; the reviewer confirms the hash independently on the object received before
  reviewing content. In the SIM2 procedure two files with the same title, the same date and
  the same version string could be told apart only by their hashes.
\item \textbf{Final standing is assigned, not claimed.} A submission that certifies itself
  ``ADMISSIBLE'' has no standing thereby. Standing is assigned by the re-audit.
\item \textbf{Order.} Declare $\to$ freeze $\to$ implement $\to$ execute $\to$ audit
  independently $\to$ assign standing. In the Block~II procedure (MOTHER-GEA,
  25--27~September 2026) implementations repeatedly arrived before the specification;
  that is exactly where the errors occurred.
\item \textbf{Two constitutional invariants} from the same procedure: \emph{solver reports
  convergence $\neq$ stationary point} (an absolute solver tolerance is not a stationarity
  test; what is needed are scale-normalised residuals evaluated independently after the
  solve) and \emph{changed frozen input $\neq$ same object}. A third finding was added: the
  same frozen inputs with two different potential functions yield two different objects;
  the specification must fix the function itself, not only its parameters.
\item \textbf{Separation of roles.} In the CTA-FUSION declaration the author declared his
  complete exposure to his own corpus and thereby excluded himself from independent review
  of its fidelity. In the SIM2 procedure (Krüger) custodian of the blind tasks, attestor and
  builder are separate roles with declarations of interest; co-authors are not regarded as
  neutral.
\item \textbf{Sealed advance prediction.} The prediction that Krüger's OP8-Neuro operator U3
  adds no structure beyond the Markov limit was deposited with a hash before the
  computation and confirmed by the result.
\item \textbf{Localise rather than judge.} A limitation in framework~A establishes or refutes
  nothing about framework~B. An author may decline a request for deeper disclosure for
  legitimate reasons; the only consequence is that independent recoverability ends at the
  declared boundary. That localises the standing; it does not condemn.
\end{itemize}

\subsection{What has proved unsuitable}

\begin{itemize}
\item \textbf{Several language-model referees from the same model.} Different roles or prompts
  increase coverage, not independence: one model has one set of blind spots. A second model
  with shared context is an echo, not a replication.
\item \textbf{Self-confirming check blocks.} A block that compares its own constants with
  themselves cannot fail and is therefore not a check.
\item \textbf{A hash over the output instead of the source.} It shows reproducibility of the
  output, not identity of the program and not physical correctness.
\item \textbf{Asserted properties.} ``Self-managing'', ``parameter-free'', ``frozen'' are
  properties that must be shown; a single fitted constant in the description is enough to
  refute ``parameter-free''.
\end{itemize}

% ============================================================================
\section{Proposal: ten building blocks}
% ============================================================================

The following building blocks are derived from the inventory. They are not a finished
framework but a proposal; a framework of rules is itself a posit and must be agreed by
those who apply it.

\subsection*{B1 -- Declared constitution before any review}
The author discloses what is posited and what is derived, in which order the layers
follow one another, and which status each statement has. For FFGFT this means: geometry
$\to$ algebra $\to$ number $\to$ independent empirical comparison. This order belongs
before the first number in every presentation.

\subsection*{B2 -- An entry path provided by the author}
Because no outsider will spend the time on the whole corpus, the burden lies with the
author: a designated entry path that shows within about an hour what everything rests on
and where it could fail. For FFGFT these are Doc.~365 (the foundation in four layers),
Doc.~377 (the interfaces to other frameworks) and the fixed predictions, for instance
$M_W=80.420$~GeV and $\sin^2\theta_{23}\in\{4/9,\,5/9\}$ (Doc.~372).

\subsection*{B3 -- Testable objects instead of texts}
What is reviewed is a frozen object: document, script, input data, each with a hash over
the source. A check script must be able to fail; it compares against an external data set
fixed beforehand, aborts on missing identity instead of printing a substitute value, and
contains null tests. Where a function is constitutive, the function itself is fixed, not
only its parameters.

\subsection*{B4 -- Predictions fixed in advance}
What a framework risks is deposited before the comparison with data, with date and hash.
A framework that fits only afterwards has risked nothing. The size of a corpus is no
substitute for a single sharp prediction.

\subsection*{B5 -- Separation of roles and interpretive authority}
The author retains interpretive authority over his posits and their declared meaning (O0).
He has no authority over whether a declared operation is actually carried out and whether
the result follows; that is decided by review. Author, reviewer and, where needed,
custodian and attestor are separate roles; each role declares its exposure. A third
party's reading of a framework is marked as a secondary rendering.

\subsection*{B6 -- A continuing register}
Reviewer findings, author responses and corrections are recorded continuously and
permanently, not overwritten. A withdrawn statement remains visible as withdrawn.
Standing is assigned, not claimed, and every assignment is contestable.

\subsection*{B7 -- Language models as tools, not as verdicts}
The use of language models is disclosed. Their results are not adopted unchecked but
secured by assertions against computation and source. A language-model review does not
replace an independent judgement; several roles of the same model are not a majority.
Responsibility for every statement remains with the person who signs it.

\subsection*{B8 -- Mark simplification as such}
Simplified presentations are marked as such and name the full version. An objection to
the simplification is not an objection to the derivation; before a reviewer declares an
error, he consults the full version.

\subsection*{B9 -- Localise rather than judge globally}
The outcome of a review is a status for each statement (for instance confirmed, not
accepted, candidate), not an overall verdict on a framework. A boundary of recoverability
localises the standing of what lies beyond it; it does not condemn it. Similar vocabulary
does not establish a bridge between frameworks.

\subsection*{B10 -- No universal checker}
Every review presupposes an agreed normalisation, and that is itself a posit. The aim is
not one checker for all frameworks but a procedure that discloses the normalisation in
each case, separates inherited conventions (computing library, system of units) from
framework-specific assumptions, and makes both recoverable.

\subsection{Graded depth of review}

To answer the question of time in practice, review can be divided into three stages that
can be carried out one after another and by different people:
\begin{center}
\small
\begin{tabular}{@{}p{0.14\linewidth}p{0.40\linewidth}p{0.38\linewidth}@{}}
\toprule
\textbf{Stage} & \textbf{Question} & \textbf{Object and effort} \\
\midrule
1 Entry & Are posits, order and status markers declared? Are there predictions fixed in advance? &
  entry documents; about an hour \\
2 Object & Do the hashes match? Do the scripts run? Can their checks fail? Do the inputs agree with the specification? &
  frozen scripts and data; hours \\
3 Depth & Does a particular chain actually follow from the posits? &
  a single derivation with its predecessors; days \\
\bottomrule
\end{tabular}
\end{center}
Stages~1 and~2 require no expertise in the subject matter of the framework and can largely
be distributed. Stage~3 does require it and is realistic only for individual chains. A
verdict at one stage holds only for that stage.

\subsection{Origin of the building blocks}

\begin{center}
\small
\begin{tabular}{@{}p{0.11\linewidth}p{0.42\linewidth}p{0.40\linewidth}@{}}
\toprule
\textbf{Block} & \textbf{Origin} & \textbf{Usability} \\
\midrule
B1 & FFGFT status markers; F O R M (Framework) & directly \\
B2 & FFGFT entry documents & directly, effort on the author's side \\
B3 & audits (hash, scripts, Block~II) & directly \\
B4 & sealed advance prediction (OP8-Neuro) & directly \\
B5 & Dot Theory O0/O2; CTA-FUSION; SIM2 & roles can be filled only with several participants \\
B6 & Doc.~190; FAH & directly \\
B7 & experience in the IPI exchange & directly \\
B8 & experience with entry texts & requires discipline on both sides \\
B9 & CIL (scope, vocabulary trap); framework~A\,$\neq$\,B & directly \\
B10 & CVP correspondence (normalisation before checker) & as a principle, not as a procedure \\
\bottomrule
\end{tabular}
\end{center}

% ============================================================================
\section{Limits and open points}
% ============================================================================

\begin{itemize}
\item \textbf{Effort.} Declaring, freezing, hashing and keeping the register take time. The
  Block~II procedure showed that the procedure reliably makes errors visible; it also
  showed that many rounds may be needed before a reviewable object exists at all.
  Procedure cannot replace content.
\item \textbf{Neutral reviewers are scarce.} In a small circle almost everyone involved is a
  co-author or correspondent. Separation of roles is then more a disclosure of closeness
  than real independence.
\item \textbf{Recognition.} A framework of rules agreed within the IPI holds at first only
  there. Whether outsiders regard an object reviewed under these rules differently from an
  unreviewed one is open.
\item \textbf{Prejudice cannot be regulated.} The building blocks make prejudice visible (for
  instance when an error verdict is issued without reference to the full version), but do
  not remove it.
\item \textbf{Interpretive authority versus review finding.} The separation in B5 is clear in
  principle, not always in the individual case: whether a statement belongs to the posit or
  is meant to follow from it must itself be declared. Here the condition applies that scope
  must be fixed at the point of derivation.
\end{itemize}

% ============================================================================
\section{Result}
% ============================================================================

\begin{enumerate}
\item Classical peer review presupposes surveyability, a shared common stock, an
  institutional prior filter and human text production; none of these assumptions holds in
  general today.
\item In the IPI exchange, in Dot Theory governance, in CIL, F O R M and FER, and in the
  FFGFT record-keeping, tools have arisen that closed these gaps in concrete cases.
\item From them ten building blocks and a graded depth of review can be derived. Most can be
  applied directly; separation of roles and recognition depend on how many participants
  there are.
\item The framework of rules is itself a posit. It gains validity only through agreement by
  those who apply it, and like every posit it must be declared openly.
\end{enumerate}

\medskip
\noindent\textbf{Check script:} \texttt{2/python/Dok378\_Skripte/pruef\_378\_verweise.py}
(checks the cited corpus documents, figures and hashes)

\medskip
\noindent\textbf{Register (Doc.~190):} no entry (new document, no correction).

\begin{thebibliography}{99}
\bibitem{dok190} J.~Pascher, \emph{Doc.~190 --- Correction Register} (continuing, 2026).
\bibitem{dok272} J.~Pascher, \emph{Doc.~272 --- FFGFT and HLV, classified according to Dot Theory} (June 2026).
\bibitem{dok274} J.~Pascher, \emph{Doc.~274 --- FFGFT in the IPI field} (June 2026).
\bibitem{dok276} J.~Pascher, \emph{Doc.~276 --- HLV homology test} (2026).
\bibitem{dok342} J.~Pascher, \emph{Doc.~342 --- Factorisation classes in G(6)} (Sept.~2026).
\bibitem{dok351} J.~Pascher, \emph{Doc.~351 --- PDG model dependence} (Sept.~2026).
\bibitem{dok361} J.~Pascher, \emph{Doc.~361 --- Kagome lattice and FFGFT} (Sept.~2026).
\bibitem{dok365} J.~Pascher, \emph{Doc.~365 --- The foundation of FFGFT in four layers} (Sept.~2026).
\bibitem{dok372} J.~Pascher, \emph{Doc.~372 --- The matrix-element principle and its reach} (Sept.~2026).
\bibitem{dok377} J.~Pascher, \emph{Doc.~377 --- Interface overview} (Sept.~2026).
\bibitem{form} S.~Vossen (ed.), D.~Haskins, J.~D.~Leavitt, J.~Pascher, J.~Quílez Zamora,
  \emph{F O R M --- Constitutional Verification Protocol}, v1.0 (2026),
  DOI 10.5281/zenodo.22875795.
\bibitem{cil} S.~Vossen, J.~Pascher, \emph{Constitutional Interpretation Layer (CIL)}, v1.1 (2026),
  \texttt{dottheory.co.uk}.
\bibitem{ccr} S.~Vossen, \emph{Constitutional Physics --- Canonical Constitutional Record}, Vols.~I--III, v1.0 (2026).
\bibitem{uc} D.~Haskins, \emph{Universal Constraint --- Unified Overview} (2026), DOI 10.5281/zenodo.20257133.
\bibitem{sim2} M.~Krüger, \emph{HLV-LIRO SIM2 v1.3.0} (2026), DOI 10.5281/zenodo.21371481.
\bibitem{ffgft} J.~Pascher, \emph{FFGFT} (2026), DOI 10.5281/zenodo.20117635;
  \texttt{github.com/jpascher/T0-Time-Mass-Duality}.
\end{thebibliography}
